\section{User Study}\label{sec:study}

\tbd{If T6 is not run, delete this section, move a short ``deployment observations'' paragraph to Section~\ref{sec:discussion}, and list the user study as future work. Every number and procedural detail below must come from the actual study.}

\paragraph{Goal and design.} We tested whether pointing with SpatialNet works for real people on the glasses, and how quickly.
\begin{itemize}
\item \emph{Participants.} \tbd{12} participants (\tbd{age range, gender, handedness, prior AR or smart-glasses use}) took part under \tbd{ethics approval or exemption, consent procedure}.
\item \emph{Setup.} \tbd{10--15} everyday objects placed near (${<}1$\,m) and far (2--4\,m) from the participant.
\item \emph{Task.} In each trial, the experimenter named an object, the participant pointed at it, and the glasses \tbd{highlighted or spoke} the localised target.
\item \emph{Trials.} \tbd{20} trials per participant (half near, half far), in \tbd{counterbalanced} order. \tbd{Optional second condition: the same task with a cloud VLM.}
\item \emph{Logging.} Each trial logged the named object, the localised object, correctness and the time from the start of the gesture to the result.
\item \emph{Questionnaires.} After the trials, participants completed the System Usability Scale~\cite{brooke1996sus} and the raw NASA-TLX~\cite{hart1988development}, followed by a short interview.
\end{itemize}

\paragraph{Results.}
\begin{itemize}
\item \emph{Task success.} \xx\% \tbd{[Wilson CI]} overall: \xx\% for near objects and \xx\% for far ones.
\item \emph{Response time.} The median time from gesture to result was \xx\,s.
\item \emph{Usability.} The SUS score was \xx\ (benchmark 68), and the raw TLX \xx.
\item \emph{Comparison.} \tbd{With two conditions: a paired Wilcoxon signed-rank test on success and time.}
\item \emph{Interviews.} \tbd{Three or four themes with short quotes, e.g., when pointing felt natural, failures participants noticed, privacy around others.}
\end{itemize}
